% called by main.tex
%
\chapter{Reproducibilidad}
\label{sec::capitulo7}

La reproducibilidad es un requisito esencial del trabajo científico: un resultado solo
adquiere valor si un tercero puede recrearlo a partir de los mismos datos, el mismo
código y el mismo procedimiento. Por este motivo, el proyecto se ha diseñado desde el
principio para ser reproducible y trazable. Este capítulo documenta los elementos que
lo hacen posible: el repositorio de código, el entorno de ejecución, la
contenedorización, el control de la aleatoriedad, la trazabilidad de los experimentos
y la integración continua.

\section{Repositorio y estructura del proyecto}
\label{sec::repro_repo}

El código fuente completo está disponible en un repositorio público de GitHub
(\url{https://github.com/SkullSupernova/MDL_TFM}). El proyecto se organiza en módulos
diferenciados, según se muestra en la Figura~\ref{fig::repro_estructura}: el código
fuente (\texttt{src/}), la configuración (\texttt{config/}), las pruebas
(\texttt{test/}), la documentación (\texttt{docs/}) y la definición de la integración
continua (\texttt{.github/}). Por otro lado, los artefactos pesados, el conjunto de
datos, los modelos entrenados (\texttt{models/}) y los resultados de los experimentos
(\texttt{experiments/}), quedan deliberadamente fuera del control de versiones, tanto
por su tamaño como por contener datos que no deben distribuirse; en su lugar, se
referencian mediante rutas y sumas de verificación.

\begin{figure}[ht]
\centering
\begin{minipage}{0.82\textwidth}
\begin{verbatim}
MDL_TFM/
|-- src/                  codigo: modelos, API, app, entrenamiento
|-- config/               configuracion (config.yml)
|-- test/                 pruebas (pytest)
|-- docs/                 documentacion
|-- .github/              integracion continua (CI/CD)
|-- models/               checkpoints (no versionado)
|-- experiments/          resultados por run (no versionado)
|-- Dockerfile
|-- docker-compose.yml
`-- requirements.txt
\end{verbatim}
\end{minipage}
\caption{Estructura del repositorio del proyecto.}
\label{fig::repro_estructura}
\end{figure}

\section{Entorno de ejecución}
\label{sec::repro_entorno}

El proyecto se desarrolla en \textbf{Python 3.12}, y todas sus dependencias se fijan de
forma explícita en el fichero \texttt{requirements.txt}, lo que permite recrear el
entorno con exactitud. Conviene distinguir dos contextos de ejecución con necesidades
distintas:

\begin{itemize}
  \item \textbf{Entrenamiento:} se realizó con \textbf{PyTorch} sobre GPU (una NVIDIA
        RTX 4060 de 8\,GB con CUDA), aprovechando la aceleración por hardware y la
        precisión mixta automática para reducir el tiempo de cómputo.
  \item \textbf{Inferencia y despliegue:} se ejecutan sobre \textbf{CPU} de forma
        deliberada. La imagen de despliegue instala las versiones de PyTorch 2.12 y
        torchvision 0.27 para CPU, lo que maximiza la portabilidad, el sistema puede
        ejecutarse en cualquier máquina, incluso sin tarjeta gráfica, a costa de una
        latencia mayor, perfectamente asumible para el uso interactivo previsto.
\end{itemize}

\section{Contenedorización}
\label{sec::repro_docker}

Para eliminar la dependencia del entorno local, la causa más frecuente de los problemas
de reproducibilidad, el sistema se distribuye como una imagen Docker construida sobre
\texttt{python:3.12-slim}, que empaqueta el código y todas sus dependencias. Existen,
además, dos vías de despliegue complementarias: la \textbf{construcción desde el código
fuente} (descrita en el Anexo~\ref{anexo::docker-local}) y la \textbf{descarga de la
imagen ya publicada} en GitHub Container Registry (Anexo~\ref{anexo::docker-ghcr}). Esta
última se publica de forma automática mediante un flujo de trabajo de GitHub Actions, de
modo que la imagen disponible siempre corresponde al estado actual del repositorio.

\section{Control de la aleatoriedad}
\label{sec::repro_semilla}

El entrenamiento de redes neuronales incorpora numerosas fuentes de aleatoriedad, la
inicialización de los pesos, el barajado de los datos o el comportamiento de ciertas
operaciones en la GPU, que, sin control, harían que dos ejecuciones idénticas
produjeran resultados distintos. Para evitarlo, todos los experimentos emplean una
\textbf{semilla fija (42)} que inicializa de forma determinista los generadores de
números aleatorios de Python, NumPy y PyTorch. Asimismo, la carga de datos se realiza en
un único hilo (\texttt{num\_workers=0}), decisión que, además de ser más robusta en
Windows, elimina la variabilidad introducida por la concurrencia. No obstante, conviene
ser honesto sobre sus límites: un determinismo \textbf{exacto a nivel de bits} no está
garantizado entre distintas máquinas o versiones de las bibliotecas de cálculo (cuDNN),
por lo que la reproducibilidad debe entenderse en términos prácticos, resultados
equivalentes dentro del margen de la incertidumbre estadística reportada, y no como una
igualdad binaria perfecta.

\section{Trazabilidad de los experimentos}
\label{sec::repro_trazabilidad}

Más allá de la semilla, cada ejecución de entrenamiento genera una carpeta
autocontenida (\texttt{experiments/<run\_id>/}) que reúne toda la información necesaria
para auditarla: la configuración empleada, las métricas de validación y prueba, las
curvas de aprendizaje, las predicciones y un informe. La pieza central de la
trazabilidad es el fichero \texttt{manifest.json}, que registra el \textit{hash} de Git
del código, el entorno y las versiones de las bibliotecas, el hardware utilizado, el
número de parámetros del modelo y la suma de verificación SHA-256 del
\textit{checkpoint}. De este modo, cada resultado queda vinculado de forma inequívoca al
estado exacto del código que lo produjo. A su vez, el índice
\texttt{experiments/leaderboard.csv} agrega todas las ejecuciones en una única tabla, y
\texttt{leaderboard\_ci.csv} añade sus intervalos de confianza, lo que constituye la
fuente autoritativa de los resultados presentados en el Capítulo~\ref{sec::capitulo4}.

\section{Integración continua y pruebas}
\label{sec::repro_ci}

Para garantizar que el código permanece funcional a lo largo del desarrollo, el
repositorio incorpora una batería de \textbf{153 pruebas automatizadas} (con
\textit{pytest}) que cubren los módulos críticos, el procesamiento de datos, los
modelos, la evaluación, el registro de experimentos y la API. Estas pruebas se ejecutan
de forma automática mediante un flujo de \textbf{integración continua} en GitHub Actions
con cada cambio enviado al repositorio, de modo que cualquier regresión se detecta de
inmediato. La integración continua se ejecuta sobre CPU de forma intencionada, para
verificar que el sistema funciona también en el entorno de despliegue.
